\begin{abstract}
Reusable chunk caches precompute document state, but the amount of model depth
that should be prepaid remains unclear. We study this question with
\textbf{CoMem}, which stores one depth-$j$ residual per context token, retrieves
a bounded chunk set, and resumes only layers $[j{:}L)$. The main result is a
matched depth-reuse frontier: with identical selected chunks, order, mask,
examples, and adapter, $j{=}12$ reduces isolated model Read from 931.9 to
664.4\,ms ($1.403\times$) while lowering a 15-cell RULER macro from 99.19 to
96.07. A rank-32 self-distillation adapter makes this operating point usable
without updating the backbone. The one-time Write amortizes only under reuse:
measured break-even is about 8--11 queries at 32k and 26--28 at 128k, depending
on generation length and store placement. We also report a transparent negative
result: at matched online latency, raw-text replay scores 64.78 versus 53.22 for
CoMem on a held-out mixed diagnostic cohort. Finally, a paired multikey
factorization bounds one failure mode. Restoring lower-layer document context
raises 92.5 to 100.0, whereas changing positions alone lowers accuracy; a
deployable 32-token overlap recovers to 98.5 without increasing persistent bytes
or per-query Read work. These findings establish a measured
quality--latency--storage frontier and a narrow mechanism diagnosis, not
superiority over raw-text retrieval or existing modular-cache systems.
\end{abstract}
